\documentclass[12pt,a4paper]{article}
\usepackage[a4paper,
    left=2.5cm,
    right=2.5cm,
    top=3.0cm,
    bottom=3.0cm,
    includehead,
    includefoot,
    headheight=60pt,
    headsep=25pt,
    footskip=30pt]{geometry}

\usepackage{setspace}
\usepackage{graphicx}
\usepackage{tabularx}
\usepackage{fancyhdr}
\usepackage{helvet}
\usepackage{xcolor}
\usepackage{titlesec}
\usepackage{hyperref}
\usepackage[brazil]{babel}

\renewcommand{\familydefault}{\sfdefault}

\renewcommand{\arraystretch}{1.3}
\setlength{\parindent}{0pt}
\setstretch{1.15}
\setlength{\headheight}{100pt}
\setlength{\headsep}{25pt}
\setlength{\footskip}{25pt}

% Caminhos das imagens --- aqui vão todas as imagens
\newcommand{\logounipar}{images/unipar_logo.png}

% --- Comando seguro: usa a imagem se existir; senão, mostra uma caixa indicando o arquivo.
% Assim o documento COMPILA no Overleaf mesmo antes de você subir os prints.
\providecommand{\imgsafe}[2][width=0.8\textwidth]{%
  \IfFileExists{#2}{\includegraphics[#1]{#2}}{%
    \framebox[0.8\textwidth]{\parbox[c][3.2cm][c]{0.75\textwidth}{\centering
      \ttfamily \small Adicione a imagem:\\ #2}}}}

% Cabeçalho institucional
\pagestyle{fancy}
\fancyhf{}
\fancyhead[L,C,R]{}
\fancyhead[C]{
\setlength{\fboxrule}{0.4pt}
\fbox{%
\begin{minipage}{0.98\textwidth}
    \begin{minipage}[c]{0.23\textwidth}
        \centering
        \vspace{2pt}
        \IfFileExists{\logounipar}{\includegraphics[width=\linewidth]{\logounipar}}{\fbox{\rule{0pt}{1.4cm}\ UNIPAR\ }}
    \end{minipage}
    \hspace{0.6cm}
    \hspace{0.7cm}
    \begin{minipage}[c]{0.65\textwidth}
        \centering
        {\normalsize\textbf{ATIVIDADE CURRICULAR DE EXTENSÃO - ACEX}}\\[5pt]
        {\small\textbf{ANÁLISE E DESENVOLVIMENTO DE SISTEMAS - ADS}}\\[-1pt]
        {\small\textbf{UNIPAR TOLEDO - 2026}}
    \end{minipage}
\end{minipage}
}\\[6pt]
\textbf{DOCUMENTO PARA REGISTRO DE ATIVIDADE DE EXTENSÃO}\\[4pt]
}

\begin{document}

\vspace*{0.2cm plus 0.1cm minus 0.9cm}

% ATENÇÃO: Os campos destacados em vermelho devem ser preenchidos conforme o projeto.
% ATENÇÃO: Remova a tag de cor (\textcolor{red}{...}) ao editar o campo.
% ATENÇÃO: Esta tabela deve permanecer em apenas uma página; resuma os conteúdos para caberem.

\begin{tabularx}{\textwidth}{|p{5cm}|X|}
\hline
\textbf{Projeto:} & \textcolor{red}{MUNDO TECH 2026} \\
\hline
\textbf{Validar:} & \textcolor{red}{000000 EXTENSÃO I} \\
\hline
\textbf{Público Atendido:} & ( ) Interno \quad (X) Externo \\
\hline
\textbf{Ação Desenvolvida:} & Letra H - Implementação de Projeto Específico de Software \\
\hline
\textbf{Integrantes:} &
60000263 - Fellype Gabriel P. dos Santos - 3º

60301340 - Pedro Kayser dos Santos - 3º \\
\hline
\textbf{Título da Atividade:} & Desenvolvimento do Sistema ConectaTech – Vitrine Digital para Loja de Informática e Assistência Técnica \\
\hline
\textbf{Resumo da Atividade:} & Foi desenvolvido o sistema ``ConectaTech'', uma página web responsiva que funciona como vitrine digital para um pequeno comércio local de informática e assistência técnica. O sistema permite ao cliente navegar pelo catálogo de produtos, buscar e filtrar itens por categoria, montar um carrinho de compras (salvo no navegador) e finalizar o pedido diretamente pelo WhatsApp, além de consultar os serviços oferecidos e enviar pedidos de orçamento. \\
\hline
\textbf{Local ou Demandante:} & Nome do comércio parceiro / cidade -- em parceria com o curso de Análise e Desenvolvimento de Sistemas, Universidade Paranaense, Toledo Campus I. \\
\hline
\textbf{Objetivos:} & - Desenvolver uma página web responsiva que sirva de vitrine digital para um pequeno comércio local.
- Aplicar conhecimentos de HTML5, CSS3 e JavaScript.
- Aproximar a tecnologia da comunidade, facilitando o acesso do público aos produtos e serviços do parceiro por meio de catálogo on-line e contato direto via WhatsApp. \\
\hline
\textbf{Etapas:} & Etapa 1 -- Levantamento de requisitos com o parceiro \textcolor{red}{(Mês/2026)}.
Etapa 2 -- Design e prototipação da interface \textcolor{red}{(Mês/2026)}.
Etapa 3 -- Desenvolvimento do catálogo, carrinho e serviços \textcolor{red}{(Mês/2026)}.
Etapa 4 -- Testes, ajustes responsivos e entrega \textcolor{red}{(Mês/2026)}. \\
\hline
\textbf{Resultados Alcançados:} & - Site funcional e responsivo finalizado.
- Catálogo com 12 produtos, carrinho persistente e pedido via WhatsApp.
- Código-fonte publicado no GitHub: \url{https://github.com/FellypePires/ConectaTech}.
- Publicação on-line (GitHub Pages): \textcolor{red}{habilitar em Settings -> Pages}. \\
\hline
\textbf{Quantidade de Público Atingido:} & Aproximadamente XX pessoas (clientes do comércio parceiro) + disponibilização da página em ambiente público. \\
\hline
\end{tabularx}

\newpage

\section*{Apresentação das Evidências}

\subsection*{Levantamento de Requisitos}

O sistema foi planejado a partir das necessidades de um pequeno comércio local de
informática que não possuía presença digital. As funcionalidades definidas foram:

\begin{itemize}
    \item Exibir um catálogo de produtos com nome, descrição, categoria e preço.
    \item Buscar produtos por nome ou descrição.
    \item Filtrar produtos por categoria (computadores, periféricos, componentes e acessórios).
    \item Adicionar produtos a um carrinho de compras, alterando a quantidade ou removendo itens.
    \item Calcular automaticamente o total do pedido.
    \item Persistir o carrinho no navegador (o pedido não se perde ao recarregar a página).
    \item Finalizar o pedido enviando um resumo automático pelo WhatsApp.
    \item Apresentar os serviços de assistência técnica oferecidos.
    \item Disponibilizar um formulário de contato/orçamento com validação dos campos.
    \item Funcionar de forma responsiva em computadores, tablets e celulares.
\end{itemize}

Tecnologias utilizadas:

\begin{itemize}
    \item HTML5 (estrutura semântica da página).
    \item CSS3 (estilização, tema e layout responsivo).
    \item JavaScript (ES6) -- lógica do catálogo, carrinho e validações.
    \item Web Storage API (\texttt{localStorage}) para persistência do carrinho.
    \item API de mensagens do WhatsApp (\texttt{wa.me}) para o envio dos pedidos.
\end{itemize}

\subsection*{Arquitetura e Estrutura do Sistema}

Por se tratar de uma aplicação web estática (executada inteiramente no navegador,
sem servidor ou banco de dados relacional), o projeto é organizado em três camadas
de arquivos. Os dados dos produtos ficam em uma estrutura em JavaScript e o carrinho
do usuário é gravado no \texttt{localStorage} do próprio navegador.

\begin{verbatim}
Extensao 1/
|-- index.html      # estrutura da página (seções e conteúdo)
|-- css/
|   |-- style.css   # estilos, tema e responsividade
|-- js/
|   |-- script.js   # lógica: catálogo, carrinho e WhatsApp
\end{verbatim}

Cada produto é representado por um objeto, conforme o exemplo abaixo:

\begin{verbatim}
{ id: 1, nome: 'Notebook i5 8GB', cat: 'computadores',
  emoji: 'X', preco: 2499.00,
  desc: 'Intel Core i5, 8GB RAM, SSD 256GB.' }
\end{verbatim}

Na Figura \ref{fig:fluxo} é apresentado o fluxo de uso do sistema, da escolha do
produto até o envio do pedido pelo WhatsApp.

\begin{figure}[h!]
    \centering
    \imgsafe[width=0.85\textwidth]{images/fluxo.png}
    \caption{Fluxo de uso: navegação no catálogo, carrinho e finalização via WhatsApp.}
    \label{fig:fluxo}
\end{figure}

\subsection*{Funções do Sistema}

A seguir, algumas das principais funções implementadas em JavaScript.

Adicionar um produto ao carrinho (cria o item ou soma a quantidade):
\begin{verbatim}
function adicionarAoCarrinho(id) {
    const produto = PRODUTOS.find((p) => p.id === id);
    const item = carrinho.find((i) => i.id === id);
    if (item) {
        item.qtd += 1;
    } else {
        carrinho.push({ ...produto, qtd: 1 });
    }
    salvarCarrinho();
    renderCarrinho();
}
\end{verbatim}

Montar a mensagem do pedido e abrir o WhatsApp:
\begin{verbatim}
function finalizarPedido() {
    let msg = 'Pedido - ConectaTech%0A%0A';
    carrinho.forEach((i) => {
        msg += '- ' + i.qtd + 'x ' + i.nome + '%0A';
    });
    msg += '%0ATotal: ' + formatarPreco(totalCarrinho());
    window.open('https://wa.me/' + WHATSAPP_NUMERO + '?text=' + msg);
}
\end{verbatim}

\subsection*{Telas do Sistema}

Na Figura \ref{fig:home} é possível visualizar a página inicial do sistema ConectaTech.

\begin{figure}[h!]
    \centering
    \imgsafe[width=0.85\textwidth]{images/home.png}
    \caption{Página inicial (\textit{home}) do sistema ConectaTech.}
    \label{fig:home}
\end{figure}

Na Figura \ref{fig:catalogo} é apresentado o catálogo de produtos com busca e filtros.

\begin{figure}[h!]
    \centering
    \imgsafe[width=0.85\textwidth]{images/catalogo.png}
    \caption{Catálogo de produtos com busca e filtro por categoria.}
    \label{fig:catalogo}
\end{figure}

Na Figura \ref{fig:carrinho} é exibido o carrinho de compras com o resumo do pedido.

\begin{figure}[h!]
    \centering
    \imgsafe[width=0.85\textwidth]{images/carrinho.png}
    \caption{Carrinho de compras com itens, quantidades e total.}
    \label{fig:carrinho}
\end{figure}

\end{document}
